\documentclass[11pt,a4paper]{article}
\usepackage[utf8]{inputenc}
\usepackage[T1]{fontenc}
\usepackage[spanish]{babel}
\usepackage{geometry}
\usepackage{enumitem}
\usepackage{titlesec}
\usepackage{xcolor}
\usepackage{hyperref}
\usepackage{fancyhdr}
\usepackage{tabularx}
\usepackage{booktabs}
\usepackage{multicol}
\usepackage{graphicx}

\geometry{margin=2.5cm}

\definecolor{musgris}{HTML}{6b7280}

\hypersetup{
    colorlinks=true,
    linkcolor=black,
    urlcolor=black,
}

\pagestyle{fancy}
\fancyhf{}
\fancyhead[L]{\textcolor{musgris}{\small MUSDIAL 2026}}
\fancyhead[R]{\textcolor{musgris}{\small Reglamento oficial}}
\fancyfoot[C]{\textcolor{musgris}{\small \thepage}}
\renewcommand{\headrulewidth}{0.5pt}

\titleformat{\section}{\Large\bfseries}{}{0em}{}[\vspace{-0.5em}\rule{\textwidth}{0.5pt}]
\titleformat{\subsection}{\large\bfseries}{}{0em}{}

\begin{document}

\begin{center}
\vspace*{2cm}
{\Huge\bfseries MUSDIAL 2026}\\[0.5cm]
{\Large Reglamento oficial del torneo}\\[2cm]
{\large Torneo de mus con formato mundial}\\[0.3cm]
{\large 25 parejas $\cdot$ 5 grupos $\cdot$ Fase de grupos + Eliminatorias}\\[2cm]
{\large\textbf{Abril -- Junio 2026}}\\[1cm]
{\normalsize Web del torneo: \url{https://musdial.softyourproblems.com}}\\[3cm]
\end{center}
\thispagestyle{empty}
\newpage

\tableofcontents
\newpage

% ============================================================
\section{Información general}
% ============================================================

El \textbf{Musdial 2026} es una competición con formato mundial de mus que se desarrollará durante tres meses: abril, mayo y junio.

\begin{itemize}
    \item \textbf{Inicio del torneo:} 1 de abril de 2026
    \item \textbf{Fecha límite de inscripción:} 30 de marzo de 2026
    \item \textbf{Inscripción:} 100\euro{} por pareja
    \item \textbf{Web del torneo:} \url{https://musdial.softyourproblems.com}
\end{itemize}

\subsection{Premios}

La recaudación total se destinará íntegramente a premios:

\begin{itemize}
    \item \textbf{1\textsuperscript{er} puesto:} 60\%
    \item \textbf{2\textsuperscript{o} puesto:} 30\%
    \item \textbf{3\textsuperscript{er} puesto:} 5\%
\end{itemize}

% ============================================================
\section{Formato de la competición}
% ============================================================

\subsection{Fase de grupos}

\begin{itemize}
    \item 5 grupos compuestos por \textbf{5 parejas} cada uno (25 parejas en total).
    \item \textbf{Duración:} abril y mayo.
    \item Cada pareja deberá disputar todos los enfrentamientos de su grupo dentro de ese plazo, siguiendo un orden de jornadas lo mejor posible.
    \item Cada enfrentamiento acordará libremente \textbf{lugar y hora} entre las parejas.
\end{itemize}

\subsection{Calendario de jornadas}

Las partidas se organizan en 5 jornadas, con 2 partidas por grupo por jornada:

\begin{center}
\begin{tabularx}{\textwidth}{lXc}
\toprule
\textbf{Jornada} & \textbf{Fechas} & \textbf{Partidas} \\
\midrule
Jornada 1 & 1 de abril -- 5 de abril & 10 \\
Jornada 2 & 6 de abril -- 19 de abril & 10 \\
Jornada 3 & 20 de abril -- 3 de mayo & 10 \\
Jornada 4 & 4 de mayo -- 17 de mayo & 10 \\
Jornada 5 & 18 de mayo -- 31 de mayo & 10 \\
\midrule
\textbf{Total fase de grupos} & & \textbf{50} \\
\bottomrule
\end{tabularx}
\end{center}

\subsection{Fase eliminatoria}

Se clasifican:
\begin{itemize}
    \item Los \textbf{3 primeros} de cada grupo.
    \item El \textbf{mejor cuarto} clasificado.
\end{itemize}

Total: \textbf{16 parejas} pasan a la fase eliminatoria.

\begin{itemize}
    \item \textbf{Octavos de final:} del 1 al 7 de junio.
    \item \textbf{Cuartos de final:} del 8 al 14 de junio.
    \item \textbf{Semifinal y final:} se jugarán en un único día, en lugar y horario por determinar. Fecha de referencia: fin de semana del 19--21 o 26--28 de junio.
\end{itemize}

\subsection{Obligaciones de las parejas}

\begin{itemize}
    \item Todas las parejas \textbf{deberán disputar la totalidad de los partidos} correspondientes a la fase de grupos.
    \item Si ambas partes no cuadran para quedar y jugar la partida, se debe avisar a la organización mediante pruebas de conversaciones. La organización pondrá un día para disputar la partida. Si una pareja no se presenta, se dará un \textbf{4--0 a la pareja que se haya presentado}.
\end{itemize}

% ============================================================
\section{Sistema de puntuación}
% ============================================================

\begin{itemize}
    \item \textbf{Victoria:} 2 puntos
    \item \textbf{Derrota:} 0 puntos
\end{itemize}

\subsection{Criterios de desempate}

\begin{enumerate}
    \item \textbf{Enfrentamiento directo} entre las parejas empatadas a puntos.
    \item \textbf{Diferencia de juegos} (juegos ganados menos juegos perdidos).
    \item En el caso del mejor 4\textsuperscript{o}, si hay empate se disputará una partida para saber quién se clasifica. Si hay triple empate, se mirará la diferencia de juegos.
\end{enumerate}

% ============================================================
\section{Normas de juego}
% ============================================================

Todas las partidas se disputarán bajo las mismas condiciones.

\subsection{Formato según fase}

\begin{itemize}
    \item \textbf{Fase de grupos:} gana la primera pareja en alcanzar \textbf{4 juegos}.
    \item \textbf{Desde octavos en adelante:} gana la primera pareja en alcanzar \textbf{6 juegos}.
    \item Cada juego se jugará a \textbf{40 tantos} (piedras).
\end{itemize}

\subsection{Reglas específicas del mus}

\begin{itemize}
    \item Se aplicará la \textbf{regla de real con figura} (el 3 no es válido).
    \item \textbf{No} se permite el perete.
    \item \textbf{No} se permiten los dúplex reales.
    \item \textbf{No} se permiten los dejes.
    \item Mus visto \textbf{se vuelve a repartir}.
    \item En pares y juego, en caso de decir no por error y luego sí, se cantará el par o el juego que lleves.
    \item Si se reparten sin querer 5 cartas, se vuelve a repartir. Y si una vez cortado el mus tiene 5 cartas, se tiran las cartas al montón jugando únicamente con las cartas de su compañero.
    \item \textbf{No} está permitido levantarse de la mesa en mitad de la partida. Solo está permitido ir al baño una vez terminada la mano o una llamada urgente. En caso de salir a fumar, se deberá acordar con los cuatro jugadores antes de empezar la partida.
\end{itemize}

% ============================================================
\section{Composición de los grupos}
% ============================================================

\begin{multicols}{2}

\textbf{Grupo A}
\begin{enumerate}[nosep]
    \item Bambel--Mora
    \item Richi--Cuchu
    \item Nule--Benito
    \item Isma--Juanillo
    \item Los nenes
\end{enumerate}

\textbf{Grupo B}
\begin{enumerate}[nosep]
    \item Felipe--Eusebio
    \item Sece--Guti
    \item Joselu--Seto
    \item Isay--Trigo
    \item Parra--Laura
\end{enumerate}

\textbf{Grupo C}
\begin{enumerate}[nosep]
    \item Chimba--Mora
    \item Calorro--Higinio
    \item Domingo--J.Miguel
    \item Pichón--J.Carlos
    \item Sonia--Kayo
\end{enumerate}

\columnbreak

\textbf{Grupo D}
\begin{enumerate}[nosep]
    \item Saúl--Diego
    \item Carlos--Use
    \item Luis--Ángel
    \item Rubén--Coco
    \item Félix--Peris
\end{enumerate}

\textbf{Grupo E}
\begin{enumerate}[nosep]
    \item Fano--Prilly
    \item Vindel--Pastor
    \item Cañamón--Cañamín
    \item Jaro--Toñín
    \item Isra--Enrique
\end{enumerate}

\end{multicols}

% ============================================================
\section{La web del torneo}
% ============================================================

El torneo cuenta con una página web donde se puede seguir todo en directo:

\begin{center}
\large\url{https://musdial.softyourproblems.com}
\end{center}

\subsection{¿Qué se puede ver en la web?}

\begin{itemize}
    \item La \textbf{clasificación} de cada grupo en tiempo real (puntos, partidas ganadas/perdidas, diferencia de juegos).
    \item Las \textbf{jornadas} con sus fechas y el estado de cada partida.
    \item El \textbf{detalle de cada partida} juego a juego, con las piedras de cada equipo.
    \item El \textbf{cuadro de eliminatorias} cuando llegue la fase.
\end{itemize}

\subsection{Enlace privado de cada pareja}

Cada pareja recibirá por WhatsApp un \textbf{enlace único y privado}. Con ese enlace se accede al panel de la pareja, desde donde se gestionan las partidas. No hace falta instalar nada, funciona directamente en el navegador del móvil.

% ============================================================
\section{Cómo usar la web (para las parejas)}
% ============================================================

A continuación se explica paso a paso cómo funciona la web. Es muy sencillo, pero es importante que las dos parejas entiendan el proceso porque \textbf{siempre hay una pareja que hace la acción y otra que la confirma}. Nunca las dos a la vez.

\subsection{Paso 1: Acceder a tu panel}

Abre el enlace que te ha mandado la organización por WhatsApp. Verás tu nombre, tu rival actual y el estado de tu partida.

\subsection{Paso 2: Empezar la partida}

Cuando las dos parejas estéis juntas y listas para jugar, \textbf{solo una de las dos} tiene que pulsar el botón. La otra espera.

\begin{enumerate}
    \item \textbf{Pareja A} abre su enlace y pulsa \textbf{``Empezar partida''}.
    \item \textbf{Pareja A} ve el mensaje: \textit{``Esperando a que tu rival confirme el inicio''}. No tiene que hacer nada más, solo esperar.
    \item \textbf{Pareja B} abre su enlace. Le aparecerá el mensaje: \textit{``Tu rival quiere empezar la partida''} y un botón \textbf{``Confirmar inicio''}.
    \item \textbf{Pareja B} pulsa \textbf{``Confirmar inicio''}.
    \item Listo. La partida está en curso para las dos parejas.
\end{enumerate}

\textbf{Importante:} No tienen que pulsar las dos parejas a la vez. Una pulsa ``Empezar'' y la otra confirma. Da igual cuál de las dos lo haga primero.

\subsection{Paso 3: Subir resultado de cada juego}

Después de jugar cada juego de mus en la mesa, hay que subir el resultado a la web. \textbf{Solo una pareja sube el resultado}, la otra lo confirma.

\begin{enumerate}
    \item \textbf{La pareja que quiera} (cualquiera de las dos, da igual cuál) pulsa \textbf{``Subir juego''}.
    \item Se le pregunta: \textbf{¿Quién ha ganado este juego?} Pulsa el nombre de la pareja ganadora.
    \item Introduce las \textbf{piedras del perdedor} (un número entre 0 y 39). Las piedras del ganador son siempre 40, no hace falta ponerlas.
    \item Pulsa \textbf{``Enviar resultado''}.
    \item La pareja que ha subido el resultado ve el mensaje: \textit{``Esperando a que tu rival confirme el juego''}. El botón de subir juego desaparece hasta que el rival confirme.
\end{enumerate}

\textbf{Importante:} No se puede subir otro juego hasta que el rival haya confirmado el anterior.

\subsection{Paso 4: Confirmar el resultado}

La otra pareja (la que NO ha subido el resultado) verá en su panel una tarjeta con el resultado y dos botones:

\begin{itemize}
    \item \textbf{``Confirmar''}: si el resultado es correcto. El juego queda registrado.
    \item \textbf{``Rechazar''}: si el resultado NO es correcto. Se avisará a la organización para que medie.
\end{itemize}

Una vez confirmado, ambas parejas pueden ver el juego en el marcador y se puede subir el siguiente juego.

\textbf{No hace falta refrescar la página.} La pantalla se actualiza sola cada pocos segundos. Cuando el rival confirme, te aparecerá automáticamente.

\subsection{Paso 5: Fin de la partida}

La partida se cierra \textbf{automáticamente} cuando una pareja alcanza:
\begin{itemize}
    \item \textbf{4 juegos ganados} en la fase de grupos.
    \item \textbf{6 juegos ganados} en la fase eliminatoria.
\end{itemize}

No hay que hacer nada más. El resultado aparece en la clasificación del grupo de forma inmediata.

\subsection{Resumen rápido del flujo}

\begin{center}
\begin{tabularx}{\textwidth}{lXX}
\toprule
\textbf{Acción} & \textbf{Pareja que actúa} & \textbf{Pareja que espera} \\
\midrule
Iniciar partida & Pulsa ``Empezar partida'' & Pulsa ``Confirmar inicio'' \\
Subir juego & Pulsa ``Subir juego'' y mete el resultado & Espera a que le aparezca \\
Confirmar juego & Espera & Pulsa ``Confirmar'' o ``Rechazar'' \\
Fin de partida & \multicolumn{2}{c}{Automático al llegar a 4 (o 6) juegos} \\
\bottomrule
\end{tabularx}
\end{center}

% ============================================================
\section{Partidas libres (de prueba)}
% ============================================================

En la página principal de la web hay un botón \textbf{``Partida libre''} que permite crear partidas de prueba para familiarizarse con el sistema antes de que empiece el torneo.

\begin{itemize}
    \item Se introduce el nombre de las dos parejas.
    \item Se puede configurar a cuántas piedras y cuántos juegos se juega.
    \item Se generan dos enlaces (uno por pareja) para jugar.
    \item Estas partidas \textbf{no afectan al torneo}.
\end{itemize}

\textbf{Se recomienda que todas las parejas hagan al menos una partida de prueba} antes del inicio del torneo para aprender a usar la web.

% ============================================================
\section{Comunicación de resultados}
% ============================================================

\begin{itemize}
    \item Los resultados se suben directamente a través de la web del torneo, con el sistema de subida y confirmación explicado anteriormente.
    \item La clasificación se actualiza en \textbf{tiempo real} en la web.
    \item La organización supervisará los resultados y mediará en caso de disputas.
\end{itemize}

% ============================================================
\section{Contacto}
% ============================================================

Para cualquier duda, incidencia o problema con la web, contactar con la organización a través del grupo de WhatsApp del torneo.

\end{document}
